\documentclass[11pt,a4paper]{article}
\usepackage[utf8]{inputenc}
\usepackage[T1]{fontenc}
\usepackage{amsmath,amsfonts,amssymb}
\usepackage{graphicx}
\usepackage[hypertexnames=false]{hyperref}
\usepackage{booktabs}
\usepackage{float}
\usepackage{geometry}
\usepackage{color}
\geometry{margin=1in}
\definecolor{zenblue}{RGB}{41,121,255}
\hypersetup{colorlinks=true,linkcolor=zenblue,urlcolor=zenblue,citecolor=zenblue}

\title{\textbf{Training on Traffic Without Retaining It}\\
\large Privacy and Efficiency Are the Same Constraint for a\\
Continually Trained Router}
\author{Zen LM Research, Hanzo AI\\
\texttt{research@hanzo.ai}}
\date{2026}

\begin{document}
\maketitle

\begin{abstract}
A router that learns from production traffic needs supervision from the requests people
actually send. The obvious way to get it---store the prompts---is a privacy failure. The
usual mitigations are worse than they look: token identifiers are a lossless encoding of
the text they came from, and frozen-backbone embeddings are invertible well enough that
published attacks recover most short inputs verbatim. Both are routinely described as
anonymization. Neither is.

We argue for the stronger and simpler position: \emph{do not retain}. What training needs
from a request is a gradient, and the request is only the courier that carries it. Once
the gradient is applied, keeping the courier adds no capability and all of the risk. We
describe a substrate that computes the update online and persists only model weights and
aggregates, and a two-tier consent model separating what trains \emph{your} router from
what trains \emph{everyone's}.

The argument's force is that it costs nothing. Retention is not a capability we give up to
buy privacy: the low-dimensional features a router actually needs are cheaper to compute,
cheaper to store, and---on the evidence of the router already serving our traffic---
sufficient. Privacy and efficiency here are not a trade-off to balance but the same
constraint seen twice.

We include an audit of our own system, which fails part of this standard today: our
routing ledger persists a frozen-backbone embedding under a code comment asserting it is
``never prompt text.'' That claim is true and misleading, and we treat finding it as the
point rather than an embarrassment to omit. A privacy argument that has not been run
against its author's own code is a marketing claim.
\end{abstract}

\section{The Problem}

A model that routes requests must learn what a good routing decision looks like, and the
only honest source of that knowledge is the traffic itself: which requests were hard, which
were trivial, which needed vision, which needed audio. Synthetic difficulty labels teach a
model to imitate the heuristic that generated them. Real traffic is the ground truth, and
it belongs to the people who sent it.

This puts an ordinary product in an uncomfortable place. The data that makes the system
good is the data it has the least right to keep. The industry's answer has largely been to
keep it anyway, under a word---``anonymized''---doing more work than it can bear.

\section{Two Mitigations That Do Not Work}

\subsection{Token identifiers are the text}

Storing token IDs rather than strings is sometimes offered as a privacy measure. It is not
one. Tokenization is a bijection with a published dictionary: $\mathrm{detokenize}(\mathrm{tokenize}(x)) = x$,
exactly, for every $x$. A token sequence is the prompt in a different alphabet, and the
key is a public download. There is no threat model in which storing token IDs is
meaningfully different from storing the prompt, because it is not different from storing
the prompt.

\subsection{Embeddings invert}

The subtler mitigation is to store an embedding: a fixed-width vector from a frozen
backbone, which ``is not the text.'' This is where the reasoning usually stops, and it
should not.

An embedding is trained to preserve meaning. Preserving meaning is exactly the property
that makes it reconstructable, and the reconstruction is not hypothetical. Inversion
attacks against text embeddings recover a large fraction of short inputs \emph{verbatim}
from the vector alone~\cite{vec2text}, treating inversion as iterative refinement in
embedding space and needing no access to the encoder's weights. Later work extends the
result across encoders and to longer inputs.

The security posture of ``we only store embeddings'' is therefore roughly the security
posture of ``we only store the prompt, lightly encrypted with a key we published.'' The
vector is smaller and harder to read by eye. It is not private.

A useful way to hold the distinction: an embedding is lossy but \emph{semantically
faithful}, and faithfulness is what inverts. Low-dimensional, purpose-built features---
length, task class, modality, observed latency---are lossy in the way that matters. They
discard the content and keep a decision. Nothing inverts a scalar.

\section{The Principle}

\begin{quote}
\emph{The gradient is what training wants. The request is the courier. Do not keep the
courier.}
\end{quote}

Once an update has been applied, the request has already given the system everything it had
to give. Retaining it buys no additional learning; it only creates an asset that can be
breached, subpoenaed, correlated, or quietly repurposed by a future version of ourselves
with different judgment. Deletion policies mitigate that risk. Never writing eliminates it.
The two are not close: a retention window is a promise, and an absent column is a fact.

This reframes the design question. Rather than asking how strongly to anonymize what we
store, ask what the smallest thing is that produces the same gradient, and whether it must
be stored at all. For a router, the answer is: less than you would guess, and no.

\section{The Substrate}

\subsection{Stream, update, drop}

The event exists long enough to produce a gradient and is then discarded. What persists is
model weights and aggregate metrics---quantities that are already the output of the
training we are permitted to do.

One concession to reality: pure online SGD over an unshuffled stream trains badly, so a
short-lived buffer (minutes) collects a batch before the update. The buffer is memory, not
a table. The distinction we are drawing is not ``fast deletion'' versus ``slow deletion''
but \emph{written} versus \emph{not written}.

\subsection{Features, not representations}

The router already computes what it needs to route. If those features are low-dimensional
and purpose-built, they are safe to hold and cheap to move. If they are a backbone
embedding, they are the prompt by another name and must not be persisted (Section~2.2).
This is one seam and it should have one form, consumed by the live router and by any
successor trained from its decisions.

\subsection{Two-tier consent}

Consent is not one bit, because the question is not one question.

\begin{table}[H]
\centering
\caption{The two scopes. Conflating them is how ``we train on your data'' becomes true in
a way the user did not agree to.}
\begin{tabular}{lll}
\toprule
\textbf{Scope} & \textbf{What it trains} & \textbf{Default} \\
\midrule
Your own head   & your router, from your traffic & opt-\emph{out} \\
The shared base & everyone's router              & opt-\emph{in} \\
\bottomrule
\end{tabular}
\end{table}

An organization's own events tuning its own router is the service working as described,
and defaults to on with an explicit way off. The same events joining a fit that ships to
everyone else is a different act with a different risk, and defaults to off until asked.
The two must not read the same flag with the same predicate, or one of them is wrong.

\subsection{What is deliberately absent}

No prompt text. No token identifiers. No backbone embeddings. No cross-org join key.
Formal guarantees (DP-SGD and its relatives) are available if the shared base fit ever
needs one, but they answer a question that arises only after you have decided to retain.
Not retaining is the cheaper instrument, and it is available first.

\section{Why This Costs Nothing}

The argument would be weaker if it were a sacrifice. It is not.

\begin{table}[H]
\centering
\caption{The same choice, read twice.}
\begin{tabular}{lll}
\toprule
\textbf{Decision} & \textbf{Privacy reading} & \textbf{Efficiency reading} \\
\midrule
Features, not embeddings & does not invert          & smaller model; no backbone pass \\
Do not retain            & nothing to breach        & no store, no lifecycle, no cost \\
Stream the update        & no corpus assembled      & no offline pipeline \\
Own head first           & data stays in its scope  & personalization beats a global fit \\
\bottomrule
\end{tabular}
\end{table}

Privacy and efficiency are usually in tension because privacy is usually bolted on: you
buy it with encryption, with auditing, with retention machinery, all of which cost. Here
the private choice and the cheap choice are the same choice, because both follow from
asking what the system actually needs and declining to acquire more. The alignment is not a
happy accident. Minimality is one property, and it is legible from either side.

There is a modeling claim inside this and we state it as a claim: a router that consumes
purpose-built features may be \emph{better} than one that reads raw tokens through a large
encoder, not merely cheaper and safer. Our production router routes across every modality
we serve on such features today. Whether that generalizes to expert-granularity routing is
untested, and Section~7 says so.

\section{Discovering the Latent Space}

A routed architecture that spans model families needs a shared representation for its
experts to interoperate through---a latent space between models and modalities. Our own
architecture papers specify a projection into that space and, in doing so, assume four
things without measuring any of them: which layers to tap, which models are worth fusing,
what the shared dimension is, and that a projection suffices at all.

Continual training turns those assumptions into measurements. A system that routes across
every modality and records where performance actually comes from is an instrument, and what
it measures is precisely which models and which layers carry signal worth aligning. The
latent space is then \emph{discovered} rather than designed---and discovered from
aggregates, which is the form we are permitted to keep. That the privacy-preserving artifact
and the architecturally informative artifact are the same object is the previous section's
point appearing a third time.

\section{Audit: Where We Fail This Today}

A privacy argument not run against its author's own system is a marketing claim. Ours,
audited at the time of writing:

\begin{table}[H]
\centering
\caption{Our own compliance with the above.}
\begin{tabular}{lll}
\toprule
\textbf{Property} & \textbf{Status} & \textbf{Note} \\
\midrule
No prompt text persisted    & \textbf{Holds}  & bodies are opt-in, redacted by default \\
Two-tier consent            & \textbf{Holds}  & own head opt-out; shared base opt-in \\
Nullable reward             & \textbf{Holds}  & unscored $\neq$ zero; dismissal records nothing \\
Counterfactual logged       & \textbf{Holds}  & the shadow arm, for off-policy learning \\
No embeddings persisted     & \textbf{FAILS}  & the ledger stores a frozen-backbone embedding \\
Do not retain               & \textbf{FAILS}  & events are written, not streamed \\
Consent read consistently   & \textbf{Unclear} & one flag, two predicates (Section~4.3) \\
\bottomrule
\end{tabular}
\end{table}

The first failure is the instructive one. Our routing ledger carries a feature column whose
contract describes it as ``the optional frozen-backbone embedding\dots\ passed through to
the routing ledger for training (never prompt text).'' Every word is accurate. The
parenthetical is nonetheless the sentence a reader will rely on, and it is false comfort:
by Section~2.2, a frozen-backbone embedding is the prompt in a form that inverts. The
comment is not a lie; it is a true statement doing a false job, which is the more common
failure and the harder one to catch. We caught it by asking what the vector was, not by
reading the assurance next to it.

The remedy follows from Section~4: the column stops being an embedding, and stops being
persisted. Both, not either.

We report this because a paper claiming a trustworthy substrate, published by an
organization that had not checked, would be evidence of nothing. The check is the
contribution. The fix is scheduled work, and this paper will be wrong about our own status
the moment it lands---which is the correct direction for a paper like this to age.

\section{What Is Not Established}

\begin{itemize}
  \item The streaming substrate is specified here, not deployed. Sections 4 and 7 must be
        read together.
  \item That features suffice for expert-granularity routing is a hypothesis. It is
        supported by a production router that routes at \emph{model} granularity on
        features, which is a weaker claim than the one we need.
  \item Non-retention is an architectural guarantee, not a formal one. It resists
        breach and subpoena because the data is absent. It does not by itself bound what a
        trained model memorizes, which is a different question with different
        instruments~\cite{dpsgd}.
  \item We publish no numbers. There is no evaluation of the substrate because it does not
        exist yet.
\end{itemize}

\section{Related Work}

Embedding inversion is the load-bearing citation: Morris et al.~\cite{vec2text} show text
embeddings can be inverted to recover most short inputs exactly, which is what disqualifies
the ``we only store vectors'' position. Differentially private
training~\cite{dpsgd} bounds what a released model reveals about any single training
example---complementary to, not a substitute for, declining to retain the example.
Federated learning~\cite{fedavg} shares the instinct that the update should travel rather
than the data; our setting is simpler, since the data is already at the server that must
learn from it, and the discipline is to not write it down. Contextual bandit
logging~\cite{bandit} gives the shape of the record---context, action, reward,
propensity---and the reason to log the counterfactual arm.

Our architecture line is specified separately: MoDE~\cite{mode} for the harvested-expert
model family, and MUEN~\cite{muen} for the router line whose v2 supplies the traffic this
paper is about.

\section*{Availability}

Router, ledger, and consent implementation: the \texttt{ai} service (private). Architecture:
\url{https://github.com/hanzoai/mode}, \url{https://github.com/hanzoai/muen}.

\begin{thebibliography}{9}

\bibitem{vec2text}
J.~X.~Morris, V.~Kuleshov, V.~Shmatikov, and A.~M.~Rush.
\emph{Text Embeddings Reveal (Almost) As Much As Text}. EMNLP 2023. arXiv:2310.06816.
\url{https://arxiv.org/abs/2310.06816}

\bibitem{dpsgd}
M.~Abadi, A.~Chu, I.~Goodfellow, H.~B.~McMahan, I.~Mironov, K.~Talwar, and L.~Zhang.
\emph{Deep Learning with Differential Privacy}. CCS 2016. arXiv:1607.00133.
\url{https://arxiv.org/abs/1607.00133}

\bibitem{fedavg}
H.~B.~McMahan, E.~Moore, D.~Ramage, S.~Hampson, and B.~Agüera y Arcas.
\emph{Communication-Efficient Learning of Deep Networks from Decentralized Data}.
AISTATS 2017. arXiv:1602.05629.
\url{https://arxiv.org/abs/1602.05629}

\bibitem{bandit}
L.~Li, W.~Chu, J.~Langford, and R.~E.~Schapire.
\emph{A Contextual-Bandit Approach to Personalized News Article Recommendation}.
WWW 2010. arXiv:1003.0146.
\url{https://arxiv.org/abs/1003.0146}

\bibitem{mode}
Zen LM Research, Hanzo AI.
\emph{MoDE v3: Mixture of Diverse Experts}.
\url{https://github.com/hanzoai/mode}

\bibitem{muen}
Zen LM Research, Hanzo AI.
\emph{MUEN v3: Multimodal Mixture of Unbound Experts}.
\url{https://github.com/hanzoai/muen}

\end{thebibliography}

\end{document}
